\chapter{Trabalhos Relacionados e Análise de Lacunas}
\label{ch:fundamentos}

Este capítulo analisa os formalismos de Redes de Petri existentes em relação às
contribuições da teoria das Redes de Petri Hierárquicas de Sinais. Em vez de
inventariar detalhes de implementação, a análise centra-se nas
\textbf{propriedades formais} --- que garantias cada formalismo oferece --- e
identifica as lacunas que esta tese aborda. A Seção~\ref{sec:fundamentos-pn}
estabelece o contexto mínimo necessário. A Seção~\ref{sec:capacidades} organiza
os trabalhos existentes por dimensão de capacidade. A Seção~\ref{sec:integracao}
examina tentativas de integração metabólico-regulatória. A
Seção~\ref{sec:lacunas} sintetiza as lacunas por meio de uma matriz de
capacidades. Para uma revisão abrangente da aplicação das Redes de Petri
a sistemas biológicos, ver~\textcite{Chaouiya2007}.

\section{Fundamentos das Redes de Petri}
\label{sec:fundamentos-pn}

\subsection{Redes de Petri Clássicas}
\label{subsec:pn-classicas}

As \textbf{Redes de Petri}~\cite{Petri1962} formalizam sistemas concorrentes por
meio da 5-tupla $\rho = (P, T, F, W, M_0)$, em que $P$ denota lugares
(condições ou recursos), $T$ denota transições (eventos ou ações),
$F \subseteq (P \times T) \cup (T \times P)$ define os arcos dirigidos,
$W: F \to \mathbb{N}^+$ atribui pesos aos arcos, e $M_0: P \to \mathbb{N}$
especifica a marcação inicial. Uma transição $t$ está \textbf{habilitada} na
marcação $M$ quando $\forall p \in {}^\bullet t: M(p) \geq W(p,t)$. O disparo
consome tokens das entradas e produz tokens nas saídas de forma atômica:
$M'(p) = M(p) - W(p, t) + W(t, p)$.

As propriedades formais incluem \textbf{P-invariantes} (conservação de tokens:
$y^T \cdot M = \text{constante}$), \textbf{T-invariantes} (sequências de disparo
cíclicas), \textbf{alcançabilidade} (marcações acessíveis a partir de $M_0$) e
\textbf{vivacidade} (ausência de impasse permanente). A complexidade da análise
é EXPSPACE-completa. Para a biologia, as Redes de Petri clássicas apresentam
limitações importantes: ausência de arcos regulatórios (catálise/inibição),
dinâmicas homogêneas, falta de semântica composicional além da independência
forte, e impossibilidade de representar a identidade elementar dos tokens.

\subsection{Técnicas de Análise Formal}
\label{subsec:analise-formal}

A análise estrutural explora a matriz de incidência $N = W(T, P) - W(P, T)$ para
calcular invariantes por álgebra linear. A verificação de modelos verifica
propriedades de lógica temporal (CTL, LTL) por exploração do espaço de estados.
Os métodos simbólicos empregam BDDs (\textit{Binary Decision Diagrams}) para
verificação eficiente de modelos de grande porte. Essas técnicas aplicam-se
igualmente a Redes de Petri clássicas e estendidas.

\section{Capacidades dos Formalismos Existentes}
\label{sec:capacidades}

\subsection{Representação de Tokens e Sistemas de Tipos}
\label{subsec:tokens}

As Redes de Petri clássicas tratam os tokens como unidades abstratas e
indistinguíveis: os lugares não diferenciam espécies moleculares além da
alocação de lugares separados, abstração suficiente para modelagem de fluxo de
controle mas limitada em expressividade bioquímica.

As \textbf{Redes de Petri Coloridas}~\cite{Jensen1981} superam essa limitação
por meio de sistemas de tipos que permitem a diferenciação de tokens através de
tipos de dados (``cores''). Lugares contêm multiconjuntos de tokens tipados, como
$\{\text{glicose}:5, \text{frutose}:2\}$. Guardas (predicados booleanos) e
expressões de arco (funções) conferem expressividade computacional. As cores
representam tipos abstratos que facilitam a verificação de protocolos e a
modelagem de sistemas de comunicação onde o tipo da mensagem importa mais do que
a composição bioquímica. Nenhum desses sistemas, porém, rastreia fórmulas
elementares como \ce{C6H12O6}, o que impede a verificação estequiométrica
automatizada.

\subsection{Dinâmicas Temporais e Semântica Multi-escala}
\label{subsec:dinamicas}

As \textbf{Redes de Petri Temporizadas}~\parencite{Ramchandani1974,Merlin1974}
introduzem atrasos determinísticos ou intervalos de tempo $[t_{\min}, t_{\max}]$
como restrições temporais. As transições disparam após durações fixas ou dentro
de janelas temporais especificadas, o que viabiliza a modelagem de processos
sequenciais com durações conhecidas --- fluxos de fabricação, protocolos de
comunicação. Essa extensão não aborda dinâmicas contínuas nem mistura de escalas
temporais distintas.

As \textbf{Redes de Petri Estocásticas}~\cite{Gillespie1977,Gibson2000,Goss1998} estabelecem
equivalência com cadeias de Markov em tempo contínuo (CTMC) por meio de tempos
de disparo distribuídos exponencialmente. Essa relação fundamental habilita
análise de estado estacionário, cálculo de vazão e métricas de desempenho por
métodos markovianos estabelecidos. A hipótese exponencial impõe dinâmicas sem
memória, simplificando a análise mas restringindo a expressividade para sistemas
com memória.

As \textbf{Redes de Petri Híbridas}~\cite{David1992} introduzem lugares contínuos
com tokens de valor real e transições contínuas que disparam a taxas diferenciais
$dM/dt = v(M)$, habilitando semântica de EDO. Modelos híbridos combinam
transições dirigidas por eventos discretos com fluxos diferenciais contínuos,
modelando com êxito cinética enzimática (Michaelis-Menten, equações de Hill),
fluxos metabólicos e acoplamento multi-escala onde estados discretos de genes
coexistem com concentrações de metabólitos contínuas. As implementações
existentes empregam escalonamento homogêneo, no qual todas as transições contínuas
disparam simultaneamente enquanto as transições discretas interrompem em tempos
de evento.

\subsection{Mecanismos Regulatórios e Tipos de Arcos}
\label{subsec:regulacao}

No formalismo clássico, não existem tipos de arcos regulatórios. Catálise e
inibição requerem codificação por lugares auxiliares ou guardas externos, tornando
as interações regulatórias invisíveis na topologia da rede.

Os \textbf{arcos de teste e arcos de inibição}~\cite{Koch2011,Heiner2008}
superam essa limitação. Os arcos de teste (ou arcos de leitura) verificam a
presença de tokens sem consumi-los, modelando catálise não consuntiva. Os arcos
de inibição bloqueiam o disparo quando a marcação de um lugar excede o limiar
$M(p) \geq \theta$, modelando inibição competitiva. Ferramentas de Redes de
Petri biológicas~\cite{Chaouiya2007} (Snoopy~\cite{Koch2011}, Charlie) implementam esses tipos de
arcos, tornando a topologia regulatória explícita. No entanto, os arcos de
inibição empregam limiares inteiros constantes ($\theta = 5$) determinados
durante a construção do modelo, incapazes de representar curvas de dose-resposta
sigmoidais da regulação bioquímica.

\subsection{Concorrência e Teoria de Independência}
\label{subsec:concorrencia}

A independência forte clássica (Mazurkiewicz, 1987) define que as transições
$t_1, t_2$ são independentes quando suas vizinhanças são completamente disjuntas:
$({}^\bullet t_1 \cup t_1^\bullet) \cap ({}^\bullet t_2 \cup t_2^\bullet) = \emptyset$.
Essa condição garante a propriedade do diamante: se ambas as transições estão
habilitadas na marcação $M$, disparar em qualquer ordem alcança a mesma marcação
final $M_{12}$. A comutatividade formal justifica a execução concorrente sem
condições de corrida. A independência forte aplica-se quando as transições
operam em pools de recursos completamente separados, o que na prática representa
menos de 15\% das transições em redes metabólicas típicas.

\subsection{Organização Hierárquica}
\label{subsec:hierarquia}

As Redes de Petri Híbridas representam sistemas multi-escala por meio de
componentes contínuos (rápidos) e discretos (lentos). Propostas como a de \textcite{Liu2013} combinam cinética enzimática com expressão gênica, mas
os modeladores classificam manualmente as transições por escala temporal sem
semântica hierárquica formal que defina relações de camada ou regras de
precedência. Essa classificação manual impede a verificação automatizada de
propriedades de controle.

\subsection{Identificação das Lacunas e Motivação da Tese}
\label{subsec:lacunas-id}

A análise das extensões existentes revela seis lacunas que motivam a teoria
das Redes de Petri Hierárquicas de Sinais.

\textbf{Composição bioquímica dos tokens.} As
Redes de Petri Coloridas oferecem sistemas de tipos abstratos que distinguem
glicose de frutose, mas os tokens carecem de fórmula elementar. Nenhum formalismo
existente rastreia cadeias como \ce{C6H12O6} para habilitar verificação
estequiométrica automatizada ou verificações de conservação de átomos.
A contribuição desta tese é a introdução da função de composição
$C: P \to \mathcal{F}$ como componente formal da 13-tupla, reservando
um campo de anotação química por lugar. Na plataforma SHYpn, esse campo
é populado por consulta ao banco de dados ChEBI~\cite{ChEBI} a partir
do nome do composto; quando a resolução é bem-sucedida, o módulo de
conservação verifica o balanço de C, H, O, N e P. A cobertura é
parcial: a heterogeneidade de notações de nomes impede resolução
automática completa, e não existe interface para entrada manual de
fórmulas --- limitação reconhecida para trabalho futuro.

\textbf{Dinâmicas heterogêneas.} As Redes de Petri
Híbridas misturam componentes discretos e contínuos, mas os sistemas biológicos
requerem quatro modos distintos coexistindo em um único modelo: cinética
enzimática contínua (Michaelis-Menten), expressão gênica estocástica (bursts
transcricionais), eventos de desenvolvimento temporizados (pontos de verificação
do ciclo celular) e dinâmicas de burst (decisão pela esporulação). As
implementações híbridas existentes empregam escalonadores homogêneos sem garantias
formais de corretude. As SHPN formalizam a coexistência de transições
heterogêneas (contínuas, estocásticas, temporizadas, burst) com um escalonador
unificado cuja corretude é demonstrada formalmente.

\textbf{Habilitação consuntiva em nível de arco.}
Os arcos de inibição clássicos empregam limiares inteiros constantes, e os arcos
de teste observam sem consumir --- nenhum dos dois fornece mecanismo para que
a \emph{avaliação de habilitação deplete o sinal avaliado}. Contudo, a biologia
da decisão celular requer exatamente essa semântica: a molécula de sinal
(e.g.\ ATP) é \emph{consumida} durante a avaliação regulatória, e esse consumo
cria fronteiras de bacia que tornam a decisão irreversível. Além disso, os
limiares dos arcos de inibição são inteiros fixados na construção do modelo,
incapazes de representar a ultrassensibilidade cooperativa descrita pela equação
de Hill ($n > 1$), que determina a nitidez da chave biológica.
Os \textbf{arcos de fluxo de sinal} das SHPN formalizam essa semântica: cada arco
carrega parâmetros $(K, n)$ derivados de $\Gamma$ que determinam o limiar efetivo
$\theta_{\text{eff}}$, avaliado como predicado booleano estrito
($M(p_s) \geq \theta_{\text{eff}} + W_s$: habilitada ou desabilitada),
e o disparo consome os tokens de sinal ($M' = M - W_s$), criando a
irreversibilidade.

\textbf{Independência fraca.} A independência forte
clássica exige vizinhanças completamente disjuntas, mas os sistemas biológicos
exibem dois padrões comuns que violam essa restrição embora permaneçam
conceitualmente concorrentes. Quando uma única enzima catalisa múltiplas reações,
as transições correspondentes compartilham arcos de teste para os lugares de
enzima; quando múltiplas reações produzem o mesmo metabólito, elas compartilham
lugares de saída. Em ambos os casos, a teoria clássica classifica as transições
como dependentes apesar da ausência de conflito de recursos. Uma das contribuições
desta tese é uma teoria de independência fraca que habilita a paralelização de
aproximadamente 65\% das transições em modelos de escala genômica.

\textbf{Formalização do controle hierárquico.} Os
modeladores designam manualmente transições como ``metabolismo'' ou ``regulação'',
mas nenhum formalismo oferece algoritmos que calculem atribuições de camada a
partir da topologia. Nenhum teorema garante que a depleção de sinais de camada
inferior desabilita transições de camada superior --- restrição fundamental para
que a exaustão de energia sobrescreva programas genéticos. A teoria de controle
hierárquico das SHPN (Capítulo~\ref{ch:shpn}) fornece o algoritmo de atribuição
de camada $\lambda: P \to \mathbb{N}$, o Teorema de Aciclicidade e a garantia de
preempção.

\textbf{Ausência de um formalismo integrado.} Tentativas de
integração metabólico-regulatória (E-Cell, Virtual Cell, Liu \emph{et al.})
requerem modelos separados conectados por interfaces de software \textit{ad hoc}.
A integração ocorre no nível de software (chamadas de API, troca de dados) e não
no nível do formalismo (semântica matemática unificada). Em consequência, os
usuários não podem calcular P-invariantes abrangendo as duas camadas nem
verificar propriedades composicionais através da interface. As SHPN fornecem um
único formalismo coerente validado pelo modelo de esporulação de
\textit{B.~subtilis} com predição de limiar dentro da incerteza experimental.

\section{Modelagem Metabólico-Regulatória Integrada}
\label{sec:integracao}

\subsection{Modelos Exclusivamente Metabólicos}
\label{subsec:apenas-meta}

A \textbf{Análise de Balanço de Fluxos} (FBA) modela o metabolismo por meio da
matriz estequiométrica $S$ com a hipótese de estado estacionário $S \cdot v = 0$,
otimizando o fluxo de biomassa sujeito a restrições. A base BiGG organiza mais
de 7.000 reconstruções em escala genômica. A programação linear fornece
otimização escalável para redes com mais de 1.000 reações, mas não há regulação
gênica, cinética nem dinâmica temporal: presume-se que todas as enzimas existam em
concentrações suficientes, ignorando o controle transcricional.

\subsection{Modelos Exclusivamente Regulatórios}
\label{subsec:apenas-reg}

As \textbf{redes Booleanas} modelam genes como variáveis binárias com regras de
atualização $\text{Gene}_i(t+1) = f(\text{Gene}_j(t), \ldots)$. Os \textbf{modelos de EDO} representam dinâmicas contínuas de concentração. Ambas as abordagens tratam os metabólitos como sinais externos e não como componentes integrados, omitindo as transformações bioquímicas que os genes finalmente controlam.

\subsection{Tentativas de Integração}
\label{subsec:tentativas-int}

A abordagem de Sim{\~a}o \emph{et al.}~\cite{Simao2005} propôs uma tradução
sistemática de redes regulatórias lógicas em Redes de Petri padrão,
integrando lógica regulatória e estequiometria metabólica num formalismo
unificado --- aplicada à biossíntese do triptofano em \textit{E.~coli}
como caso de estudo.
O \textbf{E-Cell} combina EDOs, algoritmos de Gillespie e modelos baseados em
regras de forma programática. O \textbf{Virtual Cell} acrescenta resolução
espacial por meio de EDPs de reação-difusão. A abordagem de Liu \emph{et al.}
(2013) propôs um ``quadro multi-escala'' ligando metabolismo e expressão gênica.
O padrão comum a todas essas tentativas é a necessidade de modelos separados
conectados por interfaces \textit{ad hoc}: o modelo de metabolismo exporta sinais
para o modelo de regulação gênica; o modelo de regulação exporta taxas de síntese
de enzimas de volta ao modelo de metabolismo. A integração ocorre no nível de
software, e não no nível do formalismo, impossibilitando a computação de
invariantes ou propriedades composicionais ao longo da interface.

\paragraph{SBML como formato de intercâmbio, não como formalismo.}
O \textbf{SBML} (\textit{Systems Biology Markup Language})~\cite{Hucka2003}
estabeleceu-se como o substrato \emph{de facto} de troca de modelos
bioquímicos entre laboratórios e simuladores heterogêneos, suportado por
centenas de ferramentas (COPASI, CellDesigner, libSBML) e pelo repositório
público BioModels. Trata-se, contudo, de um \emph{esquema XML de
intercâmbio} --- não de um formalismo de análise: a semântica de cada
elemento depende do simulador subjacente, e a distinção entre
participação metabólica (\texttt{reactant}), modulação não-consuntiva
(\texttt{modifier}) e consumo regulatório com formação de fronteira de
bacia não tem construto nativo --- ambas as últimas duas categorias
recaem sobre \texttt{modifier} mais uma \textit{kinetic law} customizada,
delegando a semântica ao código de cada simulador.

Essa lacuna semântica não desqualifica o SBML enquanto formato de
intercâmbio --- ela define o nicho de um \emph{formalismo} que possa
importá-lo, classificar suas anotações em $F$, $F_t$ e $F_s$, atribuir
camada hierárquica $\lambda$ aos lugares de sinal e separar
explicitamente a rede-objeto biológica do plano experimental
(parâmetros, eventos, condições). É esse o papel reservado às SHPN no
restante deste trabalho: a plataforma SHYpn (Capítulo~\ref{ch:implementacao})
opera como \emph{consumidora} de SBML, mapeando modelos curados de
BioModels para a 13-tupla SHPN acrescida do grafo de parâmetros, sem
perder os metadados de compartimento e proveniência.

\section{Síntese: Matriz de Capacidades}
\label{sec:lacunas}

A Tabela~\ref{tab:matriz-capacidades} sintetiza as seis lacunas críticas
identificadas na Seção~\ref{subsec:lacunas-id}, confrontando os formalismos
existentes com as Redes de Petri Hierárquicas de Sinais.

\begin{table}[htbp]
\centering
\caption{Matriz de lacunas de capacidade para modelagem metabólico-regulatória
integrada. Célula vazia (\textcolor{red}{\ding{55}}): capacidade ausente;
célula preenchida: nível de capacidade disponível.}
\label{tab:matriz-capacidades}
\footnotesize
\begin{tabular}{@{}lcccccc@{}}
\toprule
\textbf{Formalismo} & \rotatebox{90}{\textbf{Tokens}} & \rotatebox{90}{\textbf{Dinâmica}} & \rotatebox{90}{\textbf{Habilitação}} & \rotatebox{90}{\textbf{Concorrência}} & \rotatebox{90}{\textbf{Hierarquia}} & \rotatebox{90}{\textbf{Integração}} \\
\midrule
RP Clássica          & Abstrato  & Discreta  & \textcolor{red}{\ding{55}} & Forte   & \textcolor{red}{\ding{55}} & \textcolor{red}{\ding{55}}  \\
RP Colorida          & Tipado    & Discreta  & \textcolor{red}{\ding{55}} & Forte   & \textcolor{red}{\ding{55}} & \textcolor{red}{\ding{55}}  \\
RP Estocástica       & Abstrato  & CTMC      & \textcolor{red}{\ding{55}} & Forte   & \textcolor{red}{\ding{55}} & \textcolor{red}{\ding{55}}  \\
RP Híbrida           & Real      & Cont.+Disc.& Implícita & \textcolor{red}{\ding{55}} & \textcolor{red}{\ding{55}} & Separada \\
\midrule
\textbf{SHPN}        & \textbf{Fórmulas} & \textbf{4 modos} & \textbf{Consuntiva} & \textbf{Fraca} & \textbf{\textcolor{green}{\ding{51}}} & \textbf{Unificada} \\
\bottomrule
\end{tabular}
\end{table}

As colunas da Tabela~\ref{tab:matriz-capacidades} correspondem às seis
dimensões de capacidade avaliadas. \textbf{Tokens} designa a anotação química
via ChEBI por nome de composto (cobertura parcial; sem entrada manual de
fórmulas). \textbf{Dinâmica} designa a coexistência heterogênea de modos
contínuo, estocástico, temporizado e burst. \textbf{Habilitação} refere-se à
habilitação consuntiva por arco de sinal com limiar termodinâmico
$\theta_{\text{eff}}$ derivado de $\Gamma$. \textbf{Concorrência} designa a
independência fraca que permite enzimas compartilhadas. \textbf{Hierarquia}
corresponde ao fluxo de sinal com preempção comprovada. \textbf{Integração}
designa o formalismo metabólico-regulatório unificado.

A análise demonstra que os formalismos existentes se destacam em dimensões
individuais --- as Redes de Petri Híbridas fornecem dinâmicas contínuas, as
estocásticas estabelecem equivalência CTMC, as coloridas habilitam diferenciação
de tipos ---, mas nenhum endereça simultaneamente todas as seis dimensões. As
SHPN abrangem essas lacunas por meio de um único formalismo coerente, validado
pelo modelo de esporulação de \textit{B.~subtilis} com predição do limiar de
decisão dentro da incerteza de medição experimental.
